
This thesis investigates how a procedural music generator for tonal melodies can be developed through iterative software design. The system was built in three iterations, progressing from a simple proof of concept to a more structured generator with explicit musical objects, form handling, harmonic planning, soft constraints, and a first SATB-compatible harmonization workflow. The final system can generate both single-line melodies and simple chorale-style outputs from either automatically derived or explicitly specified harmony. The results show that richer musical representations and constraint-based generation produce more coherent and singable output than the earlier versions, while the development process also suggests that artificial intelligence can be a useful aid for implementation and refactoring when used as a tool rather than as a substitute for technical and musical understanding.
